% Post-CP19 theorem; complete root and obstruction/crystal audits passed.
\subsection{The smallest wild layer cannot hide two tame characters}

\begin{lemma}[Primitive second-level valuations]
\label{lem:B-RD55-valuations}
\status{Proved and independently checked.}
Let $p\ge5$, let $K$ be complete discretely valued of mixed
characteristic $(0,p)$ with algebraically closed residue field,
$e_K=p-1$, and $\mu_p\subset K$. Let $M$ be a potentially
supersingular height-two group whose residual inertia is
irreducible over $\mathbf F_p$. Put $q=p^2$ and $N=q-1$.
Let $L/K$ be a tame extension of degree $t\mid N$ over which
$M$ has a good formal model and all its $p$-torsion is rational.
Write $W$ for the wild image on $M[p^2]$, and let $F/L$ be
its full $p^2$-division field. Then
\[
       \operatorname{Gal}(F/L)=W\subset1+pM_2(\mathbf F_p),
       \qquad |W|\in\{p^2,p^3,p^4\}.                 \tag{B-RD55.1}
\]
Here the displayed inclusion means the kernel of reduction in
$\GL_2(\mathbf Z/p^2)$, with its usual additive matrix description.
Every nonzero $p$-torsion coordinate has $v_L$-valuation
$a/N$, where $a=(p-1)t$, and every primitive $p^2$-torsion
coordinate has valuation $a/(qN)$.
If $|W|=q$, the latter coordinates have integer $v_F$-valuation
\[
                         A=\frac{t}{p+1},\qquad p\nmid A.
                                                        \tag{B-RD55.2}
\]
\end{lemma}
\begin{proof}
The first-level valuation filtration and product argument are
those of Lemma~\ref{lem:B-RD53-leading}. They give the common
nonzero first-level valuation $\alpha=a/N$.
The negative Newton polygon of the formal $p$-series is therefore
the segment from $(1,a)$ to $(q,0)$: these are exactly its
nonzero roots of positive valuation. In particular its degree-$j$
coefficient has valuation at least $(q-j)\alpha$ for $1\le j<q$.
For a nonzero first-level point $P$, subtract its coordinate
from this series. The constant term now has valuation $\alpha$;
every intermediate coefficient is strictly above the line
from $(0,\alpha)$ to $(q,0)$, since
\[
               (q-j)\alpha>(1-j/q)\alpha\quad(0<j<q).
\]
All $q$ preimages of $P$ consequently have valuation
$\alpha/q$. These are precisely all primitive second-level points.

Wild inertia is trivial on $M[p]$ by residual irreducibility,
as proved in Lemma~\ref{lem:B-RD53-leading}. Its second-level
image is therefore an elementary abelian subgroup of the
four-dimensional additive group $M_2(\mathbf F_p)$.
The remaining quotient of the finite image is tame and is
exactly its residual image: a prime-to-$p$ quotient cannot
have a nontrivial subgroup in this reduction kernel.
Since $L$ contains the residual division field and has tame
degree, adjoining it leaves precisely $W$ as division group.
There is no residue extension, so $[F:L]=|W|$ is the
ramification degree.

The primitive coordinate valuation, normalized integrally on
$F$, is
\[
             \frac{(p-1)t|W|}{(q-1)q}
                  =\frac{t|W|}{(p+1)p^2}.
\]
It is an integer. Since $p\nmid t(p+1)$, this forces
$p^2\mid |W|$, proving (B-RD55.1). If $|W|=p^2$, the formula
is (B-RD55.2); its nonzero integer $A$ divides the prime-to-$p$
integer $t$ and is therefore itself prime to $p$.
\end{proof}

\begin{proposition}[Exclusion of the minimal wild image]
\label{prop:B-RD55-minimal-wild}
\status{Proved and independently checked.}
In the normalized remaining $p^2$ comparison of
Theorem~\ref{thm:B-RD49}, assume that the common residual
inertia representation is irreducible over $\mathbf F_p$.
If the wild image on the common $p^2$-torsion has order $p^2$,
the two requested special characters are equal.
Consequently any different-character pair in this residual
case has wild image of order $p^3$ or $p^4$.

At $p=5$, this excludes order $25$ in both normalized patterns
of Corollary~\ref{cor:B-RD51-p5}; any remaining pair with
irreducible residual inertia has wild order $125$ or $625$.
There is no CM assumption in this assertion.
\end{proposition}
\begin{proof}
Assume that the characters differ. Lemma~\ref{lem:B-RD53-leading}
and Proposition~\ref{prop:B-RD53-cycle} apply. Take the compositum
$L$ of the two faithful good tame fields and the common residual
division field. It has degree $N=p^2-1$ over $K$; in particular
it contains every tame subextension of the common second-level
division field. Put $F=L(M_1[p^2])=L(M_2[p^2])$.
By hypothesis $\operatorname{Gal}(F/L)=W$ has order $q=p^2$.
Choose a uniformizer $\varpi$ of $F$. By the preceding lemma,
every primitive second-level coordinate in either model has
integer valuation $A=p-1$, while every nonzero first-level
coordinate has valuation $qA$.

Choose Honda coordinates on both special formal groups, so
their special $p$-series is $T^q$. For the actual second-level
module $R_i=M_i[p^2](F)$ put $V_i=R_i/pR_i$. Taking the leading
coefficient of $T_i(P)/\varpi^A$ induces an injective
$\mathbf F_p$-linear map
\[
                         \ell_i:V_i\longrightarrow k.
                                                        \tag{B-RD55.3}
\]
Indeed, two primitive points have the same leading coefficient
if and only if their difference has valuation larger than $A$,
which is equivalent to their difference lying in $pR_i$.
Formal addition has ordinary addition as its leading term.
The leading coefficient of $T_i(pP)/\varpi^{qA}$ is
$\ell_i(P)^q$: the degree-$q$ term of $[p]$ is the unique
lowest-valuation term at a primitive coordinate, by the Newton
polygon calculation above, and its residue coefficient is one.
The first-level leading space is a scalar multiple of
$\mathbf F_q$, by residual irreducibility and
Lemma~\ref{lem:B-RD53-leading}. Hence the same holds for
the image in (B-RD55.3). Write it as $d_i\mathbf F_q$,
with $d_i\ne0$.

Let $g\in W$ be nontrivial, and define its lower break $b_g$ by
\[
 v_F(g\varpi-\varpi)=b_g+1,\qquad
 d_g=\overline{(g\varpi-\varpi)/\varpi^{b_g+1}}\ne0.
\]
Choose a coefficient ring $\Lambda\subset\mathcal O_K$ by
\citet[Theorem~10.160.8 and proof]{Stacks032A}, and fix a
representative $[c]\in\Lambda$ for each residue $c\in k$.
Every $g$ fixes these representatives. No canonical embedding
of the Witt ring is required. For every element $x$ of valuation
$A$, its uniformizer expansion with these fixed coefficients gives
\[
 \begin{split}
 v_F(gx-x)&=A+b_g,\\
 \overline{(gx-x)/\varpi^{A+b_g}}
       &=\overline A\,d_g\,\overline{x/\varpi^A}.
 \end{split}                                            \tag{B-RD55.4}
\]
To check the possible unit error explicitly, expand
$x=\sum_{j\ge A}[a_j]\varpi^j$. The degree-$A$ term changes
by its nonzero linear binomial term because $p\nmid A$.
Every higher degree-$j$ term changes with valuation at least
$j+b_g>A+b_g$. Thus none can cancel the displayed term.
Equivalently, a unit changes with valuation at least $b_g+1$.

Write the action of $g$ on $R_i$ as $1+pB_i(g)$, where
$B_i(g)\in\operatorname{End}_{\mathbf F_p}(V_i)$.
For every primitive point its displacement is a first-level
point. Equation (B-RD55.4) says that this displacement is
never zero. Comparing valuations gives
\[
                    b_g=(q-1)A\quad(g\ne1).             \tag{B-RD55.5}
\]
Formal subtraction has the same leading coefficient as ordinary
coordinate subtraction, since their quotient is an integral
unit with residue one. Using the $p$-series leading term now gives
\[
          \ell_i(B_i(g)v)^q=\overline A\,d_g\,\ell_i(v).
                                                        \tag{B-RD55.6}
\]
Identify $V_i$ with $\mathbf F_q$ by $\ell_i(v)=d_i z$.
Since the $q$th power fixes $\mathbf F_q$, equation (B-RD55.6)
says that $B_i(g)$ is multiplication by
\[
           \lambda_i(g)=\kappa_i d_g\in\mathbf F_q,
           \qquad \kappa_i=\overline A\,d_i^{1-q}.
                                                        \tag{B-RD55.7}
\]
The map $g\mapsto B_i(g)$ is injective on $W$, so
$\lambda_i(W)=\mathbf F_q$; define $d_g=0$ when $g=1$. Thus both models
give the same $q$-element leading ramification space, multiplied
by fixed nonzero constants $\kappa_i$.

The isomorphism of the second-level modules induces the same
linear isomorphism on $V_i$ and on $pR_i$ under multiplication
by $p$. In the two leading-coordinate identifications these
are related by $z\mapsto z^q$, which is the identity on
$\mathbf F_q$: the first-level leading coefficient is normalized
by $d_i^q$. Thus the leading-coordinate convention and tame
generator are exactly those of Lemma~\ref{lem:B-RD53-leading}.
Because the special characters differ,
Lemma~\ref{lem:B-RD53-leading} says that the residual intertwiner
must be Frobenius-semilinear, of the form
$z\mapsto c z^p$ with $c\in\mathbf F_q^\times$; the linear
case would identify the actual tangent exponents and hence
the special characters.
Commuting with every $g\in W$ consequently requires
\[
                 \lambda_2(g)=\lambda_1(g)^p.
\]
But (B-RD55.7) also gives
$\lambda_2(g)=(\kappa_2/\kappa_1)\lambda_1(g)$.
Since $\lambda_1(W)=\mathbf F_q$, this would make the
nontrivial Frobenius automorphism of $\mathbf F_q/\mathbf F_p$
a scalar multiplication on the entire field. Evaluating at
one forces that scalar to be one, and an element outside
$\mathbf F_p$ gives a contradiction.

This excludes $|W|=p^2$. The preceding lemma leaves exactly
$p^3$ or $p^4$, with the stated numerical specialization at five.
\end{proof}

\begin{remark}[Remaining scope]
\label{rem:B-RD55-scope}
\status{Proved scope statement; neither remaining wild order is excluded.}
The new restriction concerns the actual wild image on the
common second-level Galois module. It applies to non-CM as
well as CM deformations. It does not identify a group solely
from its abstract image order, and it does not assert existence
of a different-character pair of either larger order.
Reducible residual representations, and the two larger wild
orders in the irreducible case, remain unresolved here.
The exact finite supplement
\path{code/reduction_descent/verify_tame_minimal_wild_layer.py}
checks $3{,}700$ valuation cases, $166$ minimal integral cases,
$75{,}984$ field elements and separately $75{,}984$ scalar
candidates, and $1{,}180$
distinct-type orientation pairs for the primes from $5$ through
$101$. The author run passed in $0.5875725139994756$ seconds.
The independent root and obstruction/crystal runs passed in
$0.6879981279998901$ and $0.6387307759941905$ seconds,
respectively, without changing the author output. All three
statements also received complete independent mathematical audits.
The calculation does not instantiate a formal group or a congruent pair.
\end{remark}
